% arara: pdflatex: { files: [latexindent]}
\section{Fine tuning}\label{sec:finetuning}
 \texttt{latexindent.pl} operates by looking for the code blocks detailed in
 \vref{tab:code-blocks}.
 \announce{2019-07-13}{details of fine tuning of code blocks} The fine tuning of the
 details of such code blocks is controlled by the \texttt{fineTuning} field, detailed in
 \cref{lst:fineTuning}.

 This field is for those that would like to peek under the bonnet/hood and make some fine
 tuning to \texttt{latexindent.pl}'s operating. \index{warning!fine tuning}
 \index{regular expressions!fine tuning} \index{regular expressions!environments}
 \index{regular expressions!ifElseFi} \index{regular expressions!commands} \index{regular
 expressions!keyEqualsValuesBracesBrackets} \index{regular
 expressions!NamedGroupingBracesBrackets} \index{regular
 expressions!UnNamedGroupingBracesBrackets} \index{regular expressions!arguments}
 \index{regular expressions!modifyLineBreaks} \index{regular expressions!lowercase alph
 a-z} \index{regular expressions!uppercase alph A-Z} \index{regular expressions!numeric
 0-9} \index{regular expressions!at least one +} \index{regular expressions!horizontal
 space \textbackslash{h}}

 \begin{warning}
  Making changes to the fine tuning may have significant consequences for your
  indentation scheme, proceed with caution!
 \end{warning}

 \begin{widepage}
  \cmhlistingsfromfile*[style=fineTuning]{../defaultSettings.yaml}[width=.95\linewidth,before=\centering,enhanced jigsaw,breakable,yaml-TCB]{\texttt{fineTuning}}{lst:fineTuning}
 \end{widepage}

 The fields given in \cref{lst:fineTuning} are all \emph{regular expressions}. This
 manual is not intended to be a tutorial on regular expressions; you might like to read,
 for example, \cite{masteringregexp} for a detailed covering of the topic.

 We make the following comments with reference to \cref{lst:fineTuning}:
 \begin{enumerate}
  \item the \texttt{environments:name} field details that the \emph{name} of an
        environment can contain:
        \begin{enumerate}
         \item \texttt{a-z} lower case letters
         \item \texttt{A-Z} upper case letters
         \item \texttt{@} the \texttt{@} 'letter'
         \item \lstinline!\*! stars
         \item \texttt{0-9} numbers
         \item \lstinline!_! underscores
         \item \lstinline!\! backslashes
        \end{enumerate}
        \index{regular expressions!at least one +}
        The \texttt{+} at the end means \emph{at least one} of the above characters.
  \item the \texttt{ifElseFi:name} field:
        \begin{enumerate}
         \item \lstinline^@?^ means that it \emph{can possibly} begin with
               \lstinline^@^
         \item followed by \texttt{if}
         \item followed by 0 or more characters from \texttt{a-z}, \texttt{A-Z} and
               \texttt{@}
         \item the \texttt{?} the end means \emph{non-greedy}, which means `stop the
               match as soon as possible'
        \end{enumerate}
  \item the \texttt{modifyLineBreaks} field refers to fine tuning settings detailed in
        \vref{sec:modifylinebreaks}. In particular:
        \begin{enumerate}
         \item \texttt{betterFullStop} is in relation to the one sentence per line routine, detailed in
               \vref{sec:onesentenceperline}
         \item \texttt{doubleBackSlash} is in relation to the \texttt{DBSStartsOnOwnLine} and
               \texttt{DBSFinishesWithLineBreak} polyswitches surrounding double backslashes, see
               \vref{subsec:dbs}
         \item \texttt{comma} is in relation to the \texttt{CommaStartsOnOwnLine} and
               \texttt{CommaFinishesWithLineBreak} polyswitches surrounding commas in optional and
               mandatory arguments; see \vref{tab:poly-switch-mapping}
        \end{enumerate}
 \end{enumerate}

 \index{warning!capture groups}
 \begin{warning}
  For the \texttt{fineTuning} feature you should usually use \emph{non}-capturing
  groups, such as \lstinline!(?:...)! and \emph{not} capturing groups, which are
  \lstinline!(...)!
 \end{warning}

\subsection{fineTuning trailing comments}
 \begin{example}
 We can tweak the \texttt{fineTuning} for how trailing comments are classified. For
 motivation, let's consider the code given in \cref{lst:finetuning4}

 \cmhlistingsfromfile{demonstrations/finetuning4.tex}{\texttt{finetuning4.tex}}{lst:finetuning4}

 We will compare the settings given in \cref{lst:href1,lst:href2}.

 \begin{cmhtcbraster}[raster columns=2,
   raster left skip=-3.5cm,
   raster right skip=-2cm,
   raster column skip=.03\linewidth]
  \cmhlistingsfromfile[style=yaml-LST]{demonstrations/href1.yaml}[width=.1\textwidth,MLB-TCB]{\texttt{href1.yaml}}{lst:href1}
  \cmhlistingsfromfile[style=yaml-LST]{demonstrations/href2.yaml}[width=.1\textwidth,MLB-TCB]{\texttt{href2.yaml}}{lst:href2}
 \end{cmhtcbraster}

 Upon running the following commands

 \begin{commandshell}
latexindent.pl -m finetuning4.tex -o=+-mod1 -l=href1
latexindent.pl -m finetuning4.tex -o=+-mod2 -l=href2
\end{commandshell}

 we receive the respective output in \cref{lst:finetuning4-mod1,lst:finetuning4-mod2}.

 \begin{widepage}
  \cmhlistingsfromfile{demonstrations/finetuning4-mod1.tex}{\texttt{finetuning4.tex} using \cref{lst:href1}}{lst:finetuning4-mod1}

  \cmhlistingsfromfile{demonstrations/finetuning4-mod2.tex}{\texttt{finetuning4.tex} using \cref{lst:href2}}{lst:finetuning4-mod2}
 \end{widepage}

 We note that in:
 \begin{itemize}
  \item \cref{lst:finetuning4-mod1} the trailing comments are assumed to be everything following
        the first comment symbol, which has meant that everything following it has been moved to
        the end of the line; this is undesirable, clearly!
  \item \cref{lst:finetuning4-mod2} has fine-tuned the trailing comment matching, and says that
        \% cannot
        be immediately preceded by the words `Handbook', `for' or `Spoken', which means that
        none of the \% symbols have been treated as trailing comments, and the output is
        desirable.
 \end{itemize}
 \end{example}

 \begin{example}
 Another approach to this situation, which does not use \texttt{fineTuning}, is to use
 \texttt{noIndentBlock} which we discussed in \vref{lst:noIndentBlock}; using the
 settings in \cref{lst:href3} and running the command

 \begin{commandshell}
latexindent.pl -m finetuning4.tex -o=+-mod3 -l=href3
\end{commandshell}

 then we receive the same output given in \cref{lst:finetuning4-mod2}.

 \cmhlistingsfromfile[style=yaml-LST]{demonstrations/href3.yaml}[MLB-TCB]{\texttt{href3.yaml}}{lst:href3}

 With reference to the \texttt{body} field in \cref{lst:href3}, we note that the
 \texttt{body} field can be interpreted as: the fewest number of zero or more characters
 that are not right braces. This is an example of character class. \index{regular
 expressions!character class demonstration}
 \end{example}

 \begin{example}
 We can use the \texttt{fineTuning} settings to tweak how \texttt{latexindent.pl} finds
 trailing comments. \announce{2023-06-01}{fine tuning trailing comments demonstration}

 We begin with the file in \cref{lst:finetuning5}

 \cmhlistingsfromfile{demonstrations/finetuning5.tex}{\texttt{finetuning5.tex}}{lst:finetuning5}

 Using the settings in \cref{lst:fine-tuning3} and running the command

 \begin{commandshell}
latexindent.pl finetuning5.tex -l=fine-tuning3.yaml
\end{commandshell}

 gives the output in \cref{lst:finetuning5-mod1}.

 \begin{cmhtcbraster}[raster column skip=.01\linewidth,
   raster left skip=-3.5cm,
   raster force size=false,
   raster column 1/.style={add to width=-.1\textwidth},
   raster column 2/.style={add to width=-.3\textwidth},
  ]
  \cmhlistingsfromfile{demonstrations/finetuning5-mod1.tex}{\texttt{finetuning5-mod1.tex}}{lst:finetuning5-mod1}
  \cmhlistingsfromfile*[style=yaml-LST]{demonstrations/fine-tuning3.yaml}[yaml-TCB,width=0.7\textwidth]{\texttt{fine-tuning3.yaml}}{lst:fine-tuning3}
 \end{cmhtcbraster}

 The settings in \cref{lst:fine-tuning3} detail that trailing comments can \emph{not} be
 followed by a single space, and then the text `end'. This means that the
 \texttt{specialBeginEnd} routine will be able to find the pattern \lstinline!% end! as
 the \texttt{end} part. The trailing comments \texttt{123} and \texttt{456} are still
 treated as trailing comments.
 \end{example}

\subsection{fineTuning environments}
 \begin{example}
 We can use the \texttt{fineTuning} settings to tweak how \texttt{latexindent.pl} finds
 environments.

 We begin with the file in \cref{lst:finetuning6}.\announce{2023-10-13}{fineTuning
 environments}

 \cmhlistingsfromfile{demonstrations/finetuning6.tex}{\texttt{finetuning6.tex}}{lst:finetuning6}

 Using the settings in \cref{lst:fine-tuning4} and running the command

 \begin{commandshell}
latexindent.pl finetuning6.tex -m -l=fine-tuning4.yaml
\end{commandshell}

 gives the output in \cref{lst:finetuning6-mod1}.

 \begin{cmhtcbraster}[raster column skip=.01\linewidth,
   raster left skip=-3.5cm,
   raster force size=false,
   raster column 1/.style={add to width=-.1\textwidth},
   raster column 2/.style={add to width=-.3\textwidth},
  ]
  \cmhlistingsfromfile{demonstrations/finetuning6-mod1.tex}{\texttt{finetuning6-mod1.tex}}{lst:finetuning6-mod1}
  \cmhlistingsfromfile*[style=yaml-LST]{demonstrations/fine-tuning4.yaml}[MLB-TCB]{\texttt{fine-tuning4.yaml}}{lst:fine-tuning4}
 \end{cmhtcbraster}

 By using the settings in \cref{lst:fine-tuning4} it means that the default poly-switch
 location of \texttt{BodyStartsOnOwnLine} for environments (denoted
 $\BodyStartsOnOwnLine$ in \cref{tab:poly-switch-mapping}) has been overwritten so that
 it is \emph{after} the \texttt{label} command.

 Referencing \cref{lst:fine-tuning4}, unless both \texttt{begin} and \texttt{end} are
 specified, then the default value of \texttt{name} will be used.
 \end{example}

\subsection{fineTuning itemsRegex}\label{subsec:finetuning:itemsRegex}
 In \vref{lst:itemsafter} we demonstrated the default indentation
 after items. You can customise the \texttt{itemsRegex} in the
 \texttt{fineTuning}, which we demonstrate next.

 \begin{example}
 We begin with the code in \cref{lst:items2} and use the settings in
 \cref{lst:ft-itemsRegex}. Upon running

 \begin{commandshell}
latexindent.pl -l ft-itemsRegex items2-o=+-mod1
\end{commandshell}

 we receive the output in \cref{lst:items2-mod1}.

 \begin{cmhtcbraster}[raster columns=3,
   raster left skip=-3.5cm,
   raster right skip=-2cm,
   raster column skip=.03\linewidth]
  \cmhlistingsfromfile*{demonstrations/items2.tex}{\texttt{items2.tex}}{lst:items2}
  \cmhlistingsfromfile*{demonstrations/ft-itemsRegex.yaml}[yaml-TCB]{\texttt{ft-itemsRegex.yaml}}{lst:ft-itemsRegex}
  \cmhlistingsfromfile*{demonstrations/items2-mod1.tex}{\texttt{items2-mod1.tex}}{lst:items2-mod1}
 \end{cmhtcbraster}
 \end{example}

\subsection{fineTuning arguments}\label{subsec:finetuning:arguments}
 The strings allowed between arguments for \texttt{commands}, \texttt{namedGroupingBracesBrackets},
 \texttt{keyEqualsValuesBracesBrackets} and \texttt{UnNamedGroupingBracesBrackets}
 are controlled by \texttt{fineTuning}. We demonstrate with an example next.
 \announce{NEW}{fineTuning controls strings allowed between arguments }

 \begin{example}
 We begin with the code in \cref{lst:beamer1} and use the settings in \cref{lst:ft-args}.
 Upon running

 \begin{commandshell}
latexindent.pl -l ft-args beamer1-o=+-mod1
\end{commandshell}

 we receive the output in \cref{lst:beamer1-mod1}.

 \begin{cmhtcbraster}[raster columns=3,
   raster left skip=-3.5cm,
   raster right skip=-2cm,
   raster column skip=.03\linewidth]
  \cmhlistingsfromfile*{demonstrations/beamer1.tex}{\texttt{beamer1.tex}}{lst:beamer1}
  \cmhlistingsfromfile*{demonstrations/ft-args.yaml}[yaml-TCB]{\texttt{ft-args.yaml}}{lst:ft-args}
  \cmhlistingsfromfile*{demonstrations/beamer1-mod1.tex}{\texttt{beamer1-mod1.tex}}{lst:beamer1-mod1}
 \end{cmhtcbraster}
 \end{example}

\subsection{fineTuning ifElseFi}\label{subsec:finetuning:ifElseFi}
 We can fine tune the \texttt{elseOrRegEx} for \texttt{ifElseFi} blocks as we demonstrate next.

 \begin{example}
 We begin with the code in \cref{lst:ft-ifelsefi} and use the settings in
 \cref{lst:ft-ifelsefi-yaml}. Upon running

 \begin{commandshell}
latexindent.pl -l ft-ifelsefi ft-ifelsefi-o=+-mod1
\end{commandshell}

 we receive the output in \cref{lst:ft-ifelsefi-mod1}.

 \begin{cmhtcbraster}[raster columns=3,
   raster left skip=-3.5cm,
   raster right skip=-2cm,
   raster column skip=.03\linewidth]
  \cmhlistingsfromfile*{demonstrations/ft-ifelsefi.tex}{\texttt{ft-ifelsefi.tex}}{lst:ft-ifelsefi}
  \cmhlistingsfromfile*{demonstrations/ft-ifelsefi.yaml}[yaml-TCB]{\texttt{ft-ifelsefi.yaml}}{lst:ft-ifelsefi-yaml}
  \cmhlistingsfromfile*{demonstrations/ft-ifelsefi-mod1.tex}{\texttt{ft-ifelsefi-mod1.tex}}{lst:ft-ifelsefi-mod1}
 \end{cmhtcbraster}
 \end{example}

\subsection{fineTuning lookForAlignDelims}\label{subsec:finetuning:lookForAlignDelims}
 As detailed in \vref{subsec:align-at-delimiters} there is a basic and an advanced way to specify the entries
 in \texttt{lookForAlignDelims}; any entries not specified explicitly, will be read
 from the default values specified in \texttt{fineTuning}.

 As an example, the settings in \cref{lst:ft-align1,lst:ft-align2,lst:ft-align3} are
 equivalent.

 \begin{cmhtcbraster}[raster columns=3,
   raster left skip=-3.5cm,
   raster right skip=-2cm,
   raster column skip=.03\linewidth]
  \cmhlistingsfromfile*{demonstrations/align1.yaml}[yaml-TCB]{\texttt{align1.yaml}}{lst:ft-align1}
  \cmhlistingsfromfile*{demonstrations/align2.yaml}[yaml-TCB]{\texttt{align2.yaml}}{lst:ft-align2}
  \cmhlistingsfromfile*{demonstrations/align3.yaml}[yaml-TCB]{\texttt{align3.yaml}}{lst:ft-align3}
 \end{cmhtcbraster}

\subsection{fineTuning using -y switch}
 \begin{example}
 You can tweak the \texttt{fineTuning} using the \texttt{-y} switch, but to be sure to
 use quotes appropriately. For example, starting with the code in \cref{lst:finetuning3}
 and running the following command

 \begin{commandshell}
latexindent.pl -m -y='modifyLineBreaks:oneSentencePerLine:manipulateSentences: 1, modifyLineBreaks:oneSentencePerLine:sentencesBeginWith:a-z: 1, fineTuning:modifyLineBreaks:betterFullStop: "(?:\.|;|:(?![a-z]))|(?:(?<!(?:(?:e\.g)|(?:i\.e)|(?:etc))))\.(?!(?:[a-z]|[A-Z]|\-|~|\,|[0-9]))"' issue-243.tex -o=+-mod1
\end{commandshell}

 gives the output shown in \cref{lst:finetuning3-mod1}.

 \cmhlistingsfromfile{demonstrations/finetuning3.tex}{\texttt{finetuning3.tex}}{lst:finetuning3}
 \cmhlistingsfromfile{demonstrations/finetuning3-mod1.tex}{\texttt{finetuning3.tex} using -y switch}{lst:finetuning3-mod1}
 \end{example}
